% ============================================================
% Project: THE PHILOSOPHY OF VIABILITY
% Book: Book I — Ontology of Viability
% Book ID: VS.B01
% Part: Part I — Identity and Manifestation
% Chapter: Chapter 07 — Formal and Logical Reconstruction
% File: 4_ASSETS/tables/VS.B01.P01.C07.S02.SS04.P001.TL01.ontological_vocabulary_and_formal_roles.tex
% Asset: TL01 — Ontological Vocabulary and Formal Roles
% Purpose: To distinguish foundational ontological terms,
% related concepts, and formal constraints governing their use.
% ============================================================

\begin{table}[htbp]
\centering
\small
\caption{Ontological Vocabulary and Its Formal Roles}
\label{tab:ch07_ontological_vocabulary}
\renewcommand{\arraystretch}{1.25}
\begin{tabular}{|p{0.19\textwidth}|p{0.29\textwidth}|p{0.42\textwidth}|}
\hline
\textbf{Term} &
\textbf{Role in the Ontological Inquiry} &
\textbf{Constraint on Formal Use}
\\
\hline
Being &
The actuality of what exists. &
Must remain analytically distinguishable from identity while ontologically inseparable from it.
\\
\hline
Existent &
That which is considered as existing and as the existent that it is. &
Must not be presupposed to be exclusively a substance, aggregate, process, or composition.
\\
\hline
Identity &
The existent as the existent that it is. &
Must not be reduced to determination, unity, numerical identity, invariance, manifestation, or continuity.
\\
\hline
Determination &
The determinate character of an existent. &
The constitutive role of a particular determination must be established rather than inferred from its presence.
\\
\hline
Distinction and Difference &
Analytical or ontological differentiation under a specified respect. &
Conceptual distinction must not be treated as proof of ontological separation; difference must be specified.
\\
\hline
Relation &
A connection or ordering among terms under a specified respect. &
The relata, domain, and intended relational role must be identified.
\\
\hline
Manifestation &
The manifest presence or presentation of an existent or its aspects. &
Must not be treated as an exhaustive substitute for the identity of the existent.
\\
\hline
Expression &
The articulation or expression of aspects associated with an existent. &
Must remain distinguishable from manifestation unless their relation is explicitly established.
\\
\hline
Unity and Plurality &
Multiplicity or unity considered under a specified respect. &
Component, manifestational, relational, and identity plurality must not be conflated.
\\
\hline
Invariance &
Preservation under a specified transformation and respect. &
Transformation, domain, respect, type, and limits must be specified.
\\
\hline
Transformation &
Change between conditions or configurations. &
Must not be taken to establish either preservation or replacement of identity by itself.
\\
\hline
Persistence and Continuity &
Continuation across an interval or sequence of conditions. &
Must be distinguished from identity and specified with respect to subject, interval, and conditions.
\\
\hline
Structure and Composition &
Organization of relations or constitution through components. &
Neither structure nor composition may be assumed to exhaust identity.
\\
\hline
Numerical Identity &
Being one and the same object within a specified logical interpretation. &
The logical equality symbol must not be treated as a complete expression of ontological identity.
\\
\hline
Signature &
A structured characterization associated with an existent. &
Descriptive or identifying usefulness does not establish ontological constitution.
\\
\hline
\end{tabular}
\end{table}

% ============================================================
% END OF ASSET TL01 — CHAPTER 07, SECTION 02
% ============================================================